% modello documentale di nicola santi ego@nicolasanti.it

\documentclass[pdftex,a4paper,11pt,twoside,openright]{book}

\usepackage[italian]{babel}
\usepackage[T1]{fontenc}
\usepackage[utf8]{inputenc}
\usepackage{indentfirst}

% decora i titoli dei capitoli
\usepackage[Lenny]{fncychap}

% inserimento di immagini come effetto indesiderato
% produce solo pdf come output e non più dvi.
\usepackage[pdftex]{graphicx}
\DeclareGraphicsExtensions{.pdf,.png,.jpg}
\graphicspath{{./images/}}

\newcommand{\HRule}{\rule{\linewidth}{0.5mm}}

% floatings, caption laterali
\usepackage{sidecap}

% immagini avvolte all'interno del testo
\usepackage{wrapfig}

% citazioni in stile Hardvare
\usepackage{natbib}
\bibpunct{(}{)}{;}{a}{,}{,}

% forza LaTeX ad una spaziatura fra parole non inglese
\frenchspacing 

% per inserire del codice sorgente
\usepackage{listings}
\usepackage{color}
\definecolor{light-gray}{gray}{0.95}


\lstnewenvironment{java}[1][]
{\lstset{language=Java,numbers=left, numberstyle=\tiny, stepnumber=2, numbersep=5pt,backgroundcolor=\color{light-gray},breaklines=true,breakatwhitespace=true,breakindent=0pt,basicstyle=\scriptsize,#1}}
{}

\lstnewenvironment{bash}[1][]
{\lstset{language=Bash,numbers=left, numberstyle=\tiny, stepnumber=2, numbersep=5pt,backgroundcolor=\color{light-gray},breaklines=true,breakatwhitespace=true,breakindent=0pt,basicstyle=\scriptsize,#1}}
{}

\lstnewenvironment{xml}[1][]
{\lstset{language=Xml,numbers=left, numberstyle=\tiny, stepnumber=2, numbersep=5pt,backgroundcolor=\color{light-gray},breaklines=true,breakatwhitespace=true,breakindent=0pt,basicstyle=\scriptsize,#1}}
{}



% indentazione del primo paragrafo
\setlength{\parindent}{18pt}

% serve per gestire le tabelle
\usepackage{array}
\usepackage{multirow}
\usepackage{booktabs}

% per inserire l'euro
\usepackage[official]{eurosym}

% liste più compatte
\usepackage{mdwlist}

% liste inline
\usepackage{paralist}

% miglioare i riferimenti
\usepackage[italian]{varioref}
%\def\reftextfaraway#1{a pagina~\pageref{#1}}%
%\def\reftextpagerange#1#2{nelle pagine~\pageref{#1}--\pageref{#2}}%

% hyper link in LaTex
\usepackage{html}

% indici
\usepackage{makeidx}
\makeindex

% un simpatico box

\usepackage{dashbox}

\newenvironment{nota}
{\rule{1ex}{1ex}\hspace{\stretch{1}}}
{\hspace{\stretch{1}}\rule{1ex}{1ex}}



\begin{document} 

%-----------TOP MATTER

%\title{Pojo's in action} 

%\author{Nicola Santi\\
%  Università di Bologna\\
%  \texttt{ego@nicolasanti.it}}
%\date{\today}

%\maketitle 
\begin{titlepage}
 
\begin{center}
 
 
% Upper part of the page
\includegraphics[width=0.40\textwidth]{regola-kit}\\[1cm]
 
\textsc{\LARGE }\\[1.5cm]
 
\textsc{\Large java ed il web per gente pratica }\\[0.5cm]
 
 % Title
\HRule \\[0.4cm]
{ \huge \bfseries scrivere con Regola}\\[0.4cm]
 
\HRule \\[1.5cm]

 
% Author and supervisor
\begin{minipage}{0.4\textwidth}
\begin{flushleft} \large
\emph{Autori:}\\
Nicola \textsc{Santi}, Marco \textsc{Rosi}
\end{flushleft}
\end{minipage}
\begin{minipage}{0.4\textwidth}
\begin{flushright} \large
\emph{Supervisore:} \\
Lorenzo \textsc{Bragaglia}
\end{flushright} 
\end{minipage}
 
\vfill
 
% Bottom of the page
{\large \today}
 
\end{center}
 
\end{titlepage}

% -----------TOC
\tableofcontents
%\listoffigures
%\listoftables


\subsection*{Prefazione}
\addcontentsline{toc}{subsection}{Prefazione}

% -----------CAPITOLI
\chapter{Getting Started}

Questo capitolo è una cura per gli impazienti; seguendo le istruzioni dei paragrafi seguenti sarete in grado di installare e predisporre una prima applicazioni web con Regola.

\section{Installare Regola kit}
\index{installazione} 
Le applicazioni realizzate con Regola kit utilizzano Maven 2 per tutta la gestione del ciclo di build (compilazione, esecuzione dei test, creazione dei file war, \ldots). Quindi come prima cosa bisogna installare Maven 2 
scaricandolo dal sito \htmladdnormallink{maven.apache.org} {http://maven.apache.org} e seguendo le istruzioni.

Finita l'installazione di Maven 2 verificate che tutto sia andato a buon fine aprendo una console di terminale e lanciando il seguente comando.

\begin{bash}
nicola@casper:~# mvn -version
Maven version: 2.0.8
Java version: 1.6.0_03
OS name: "linux" version: "2.6.22-14-generic" arch: "i386" Family: "unix"
\end{bash}

Se tutto è corretto Maven 2 risponde al comando restituendo la sua versione, quella di Java ed infine alcune informazioni circa il sistema operativo in uso.

Regola kit non richiede nessuna installazione particolare (anche se è possibile scaricare un pacchetto contente documentazione e comandi di utilità): quindi finita l'installazione di Maven 2 siete pronti già per utilizzare Regola kit.


%\setlength{\fboxrule}{0.1pt} 
%\framebox[\width]{
%\begin{minipage}{0.9\textwidth}
\begin{nota}
Per maggiori informazioni su come installare Regola kit sulle vostre macchine di sviluppo si rimanda al capitolo  \vref{chap:installazione}.
\end{nota}
%\end{minipage}
%}


\section{Predisporre il database}
Per questo tutorial ipotizziamo di avere a disposizione un database di tipo MySql già installato sulla macchina dove intendiamo creare il progetto. Assicuratevi che il database stia funzionando e digitate il comando seguente:

\begin{bash}
nicola@casper:~# mysql -u root -p
\end{bash}

Vi troverete dentro la shell di amministrazione di MySql. Approfittatene per creare un nuovo database che sarà utilizzato dall'applicazione digitando il comando.

\begin{bash}
mysql> create database clienti;
\end{bash}

(Attenzione al punto e virgola in fondo al comando).
A questo punto non ci resta che creare anche l'utente utilizzato dalla nostra applicazione per accedere al database (nell'esempio porta il mio nome \emph{nicola}).

\begin{bash}
mysql> grant all on clienti.* to 'nicola'@'localhost';
\end{bash}

Bene, con il database abbiamo finito. Digitate questo ultimo comando per uscire dalla shell di amministrazione di MySql e passare al paragrafo seguente.

\begin{bash}
mysql> exit
\end{bash}

\begin{nota}
Regola kit è in grado di utilizzare diversi DBMS (ad esempio Oracle, Microsoft Sql, PostgreSQL, Hypersonic, \ldots). Per sapere come configurare la vostra applicazione per utilizzare DBMS diversi da MySql si rimanda al capitolo \vref{chap:persistenza}.
\end{nota}



\section{Creare un progetto con Regola kit}

Posizionatevi nella cartella dove volete creare il vostro nuovo progetto e digitate un comando simile al seguente con l'accortezza però di modificare il parametro gruopId (nell'esempio \emph{com.acme}) con il nome del vostro package di default ed il parametro  artifactId (nell'esempio  \emph{clienti}) con il nome del nuovo progetto.

\begin{bash}
nicola@casper:~# mvn archetype:create -DarchetypeGroupId=org.regola  \ 
-DarchetypeArtifactId=regola-jsf-archetype \
-DarchetypeVersion=1.1-SNAPSHOT -DgroupId=com.acme \
-DartifactId=clienti
\end{bash}

\emph{Attenzione: il comando qui sopra è riparito su diverse righe per chiarezza tipografica, deve invece essere digitato su una sola riga.}

La prima volta che lanciate questo comando Maven 2 scarica tutte le librerie necessarie (la cosa potrebbe prendere un po' di tempo) e crea una sotto cartella col nome del progetto (nell'esempio la cartella clienti). Il progetto è questo punto è già stato creato, posizioniatevi all'interno della cartella  \emph{clienti} col comando:

\begin{bash}
nicola@casper:~# cd clienti
\end{bash}


\section{Collegarsi al database}\label{sec:db-link}

Prima di lanciare la nostra nuova applicazione è necessario informarla circa le coordinate del database da utilizzare, per farlo bisogna apportare un modifica al file src/test/resources/jetty/env.xml con il vostro editor di testo preferito. Dovete inserire cambiare solo il nome del database (alla riga 6) e lo username (riga 8) da utilizzare per ottenere qualcosa di simile al frammento di xml seguente:

\begin{xml}
...
<New id="jira-ds" class="org.mortbay.jetty.plus.naming.Resource">
  <Arg>jdbc/Datasource</Arg>
  <Arg>
    <New class="org.enhydra.jdbc.standard.StandardConnectionPoolDataSource">
      <Set name="Url">jdbc:mysql://localhost/clienti</Set>
      <Set name="DriverName">com.mysql.jdbc.Driver</Set>
      <Set name="User">nicola</Set>
    </New>
  </Arg>
</New>
...
\end{xml}

\begin{nota}
Per avere maggiori informazioni sulle diverse configurazioni relative alle connessioni al database si rimanda al capitolo \ref{chap:database}
\end{nota}


\section{Avviare l'applicazione}
 \index{applicazione!avviare}
Ora tutto è pronto per avviare l'applicazione, se avete lasciato la cartella principale del progetto tornateci e da lì lanciate il comando seguente (e lasciate aperta la console):

\begin{bash}
nicola@casper:~/projects/clienti# mvn jetty:run
\end{bash}

Sullo schermo si susseguiranno diverse righe per informarvi che l'applicazione è stata inizializzata e quando, infine, apparirà la dicitura \emph{Started Jetty Server} saprete che tutto è pronto.

Lasciando sempre aperta la console aprite un'istanza del vostro browser e collegatevi all'indirizzo \htmladdnormallink{localhost:8080/clienti} {http://localhost:8080/clienti} per vedere la pagina di benvenuto della vostra applicazione.

Complimenti, avete appena fatto il primo passo nel mondo delle applicazioni Regola kit!

\begin{nota}
Per visualizzare l'applicazione state utilizzando un piccolo (ma molto completo) application server di nome Jetty. Per la messa in produzione però si consiglia di utilizzare dei container diversi (ad esempio Tomcat o JBoss). Per imparare a come creare i pacchetti per questi application server si rimanda al capitolo \vref{chap:produzione} 
\end{nota}

\section{Struttura di un progetto Regola kit}
La struttura della cartella di un progetto Regola kit si impronta alla struttra standard di un progetto web di Maven 2. Al primo livello troviamo:

\begin{center}
{
  \begin{tabular}{ | l | p{9cm} | }
  \hline
  pom.xml & il file di configurazione di Maven 2 \\ \hline
  src/ & la cartella dei sorgenti  \\ \hline
  target/ & contiene i file compilati ed i pacchetti per le consegne \\ \hline
  \end{tabular}
}
\end{center}

La cartella target contiene quanto i pacchetti pronti per la consegna con la classi compilate ed i descrittori. Si tratta di una cartella il cui contenuto è ricreato ogni volta si lanci il comando:

\begin{bash}
nicola@casper:~/projects/clienti# mvn package
\end{bash}

La cartella src contiene i sorgenti (html, js, css e java) dell'applicazione. Al suo interno potete trovare:

\begin{center}
{
  \begin{tabular}{ | l | p{9cm} | }
  \hline
  main/ & contiene i sorgenti dell'appplicazione \\ \hline
  test/ & contiene i sorgenti dei test \\ \hline
  site/ & reportistica generata da Maven 2  \\ \hline
  \end{tabular}
}
\end{center}

La cartella main e la cartella test contengono entrambe i sorgenti java (nella sottocartella java) e le altre risorse (nella cartella risorse). Queste ultime sono i file di configurazione, i mappaggi orm e, in generale, tutto quello che non sono sorgenti Java ma devono finire comunque nel classpath. La differenza tra la cartella main e quella test e che il contenuto di quest'ultima non finisce mai nei pacchetti destinati alla produzione ma è usato esclusivamente per l'esecuzione dei test.
Infine la cartella main contiene anche la sottocartella webapp dove si trova la webroot, ovvero le pagine web dell'applicazione ed il file web.xml.

\begin{nota}
Per una descrizione completa dei file e delle cartelle standard di un progetto Regola kit si rimanda al capitolo \vref{chap:struttura}
\end{nota}


\section{Persistenza su database}

Regola kit abbraccia una metodologia di sviluppo incentrata sull'analisi del dominio del problema (Domain Driven Development) per cui il primo passo è quello di creare le classi di modello. Spesso queste classi sono persistite sul database per cui si potrebbe iniziare scrivendo la classe e poi creando la corrispondente tabella sul database. Oppure, al contrario come faremo tra poco, creando prima la tabella del database e facendoci poi creare in automatico la classe Java (nel prossimo paragrafo \vref{sec:gsmodello}). Colleghiamoci nuovamente al database clienti.

\begin{bash}
nicola@casper:~/projects/clienti# mysql -u root -p clienti
\end{bash}

Creiamo una piccola tabella con la sola chiave primaria (id) ed un campo di descrizione (label).

\begin{bash}
mysql> create table prodotti (id int(11) not null auto_increment, label varchar(80) not null, primary key (id) );
mysql> describe prodotti;
+-------+-------------+------+-----+---------+----------------+
| Field | Type        | Null | Key | Default | Extra          |
+-------+-------------+------+-----+---------+----------------+
| id    | int(11)     | NO   | PRI | NULL    | auto_increment | 
| label | varchar(80) | NO   |     |         |                | 
+-------+-------------+------+-----+---------+----------------+
2 rows in set (0.02 sec)
\end{bash}

Inseriamo qualche dato nella tabella, ad esempio alcune descrizioni di esempio per verificare poi il funzionamento dell'applicazione.

\begin{bash}
mysql> insert into prodotti values (null, 'book');
Query OK, 1 row affected (0.05 sec)

mysql> insert into prodotti values (null, 'bottle');
Query OK, 1 row affected (0.00 sec)

mysql> insert into prodotti values (null, 'paper');
Query OK, 1 row affected (0.00 sec)

mysql> select * from prodotti;
+----+--------+
| id | label  |
+----+--------+
|  1 | book   | 
|  2 | bottle | 
|  3 | paper  | 
+----+--------+
3 rows in set (0.01 sec)
\end{bash}

Adesso il database contiene una tabella con dei dati che possiamo usare per persistere le nostre classi di modello. Usciamo dal database e torniamo all'applicazione.

\begin{bash}
mysql> exit
\end{bash}

\section{(Ri)collegarsi al database}
Al paragrafo \vref{sec:db-link} abbiamo configurato l'applicazione per utilizzare il nostro database in fase di esecuzione. Adesso dobbiamo fare in modo che anche in fase di sviluppo si possa accedere al database (ad esempio per usare i generatori di codice o lanciare la batteria di test). Il file da modificare è src/test/resources/designtime.properties e deve essere aggiornato in modo da contenere lo username, la password ed il nome del database. Il risultato finale deve risultare simile al seguente:

\begin{bash}
...
hibernate.dialect=org.hibernate.dialect.MySQLDialect
hibernate.connection.driver_class = com.mysql.jdbc.Driver
hibernate.connection.url = jdbc:mysql://localhost/clienti
hibernate.connection.username = nicola
hibernate.connection.password = 
...
\end{bash}


\section{Classi di modello}\label{sec:gsmodello}
Siete ora in grado di scrivere la classe di modello che sarà persistita sulla tabella prodotti... oppure potete farvela generare automaticamente e poi modificare convenientemente le classi prodotte. Userete gli Hibernate Tools che sono già configurati all'interno delle applicazioni prodotte con Regola kit e trovano nell'unico file src/test/resources/hibernate.reveng.xml la configurazione di tutto il processo di generazione inversa, a partire cioè dal database. Specificate il nome della tabella da cui partire (\emph{prodotti} alla riga 2), il nome della classe (\emph{Prodotto}, al singolare e con la prima lettera maiuscola nella riga 4), il package da utilizzare (\emph{com.acme.model} sempre alla riga 2).

\begin{bash}
...  
  <table-filter match-name="prodotti"  package="com.acme.model" exclude="false"/>

   <table name="prodotti" class="Prodotto" >
      <primary-key property="id" />
   </table>
...
\end{bash}

Ora avviate la generazione: posizionantevi nella cartella principale del vostro progetto ed utilizzate il plugin di Maven 2 Hibernate3 che consente di generare le classi java (il goal hbm2java), i file di mappaggio di hibernate (hbm2hbmxml) e la configurazione generale di hibernate (hbm2cfgxml). 

\begin{bash}
nicola@casper:~/projects/clienti# mvn hibernate3:hbm2java hibernate3:hbm2hbmxml hibernate3:hbm2cfgxml
\end{bash}

Il primo file generato si trova nella posizione src/main/java/com/acme/model/Prodotto.java e contiene la classe Java:


\begin{java}
public class Prodotto  implements java.io.Serializable {
    
    private Integer id;

    public Prodotto() {
    }

    public Integer getId() {
        return this.id;
    }
    
    public void setId(Integer id) {
        this.id = id;
    }
}
\end{java}

Poi è stato creato il file con i mappaggi di hibernate src/main/resources/com/acme/model/Prodotto.hbm.xml, molto semplice in questo caso:

\begin{xml}
<?xml version="1.0"?>
<!DOCTYPE hibernate-mapping PUBLIC "-//Hibernate/Hibernate Mapping DTD 3.0//EN"
"http://hibernate.sourceforge.net/hibernate-mapping-3.0.dtd">
<!-- Generated 14-apr-2008 12.23.38 by Hibernate Tools 3.2.0.CR1 -->
<hibernate-mapping>
    <class name="com.acme.model.Prodotto" table="prodotti" catalog="clienti">
        <id name="id" type="java.lang.Integer">
            <column name="id" />
            <generator class="identity" />
        </id>
    </class>
</hibernate-mapping>
\end{xml}

Ed infine è stato inserito un riferimento a quest'ultimo file di mappaggio dentro la configuraizone principale di Hibernate    src/main/resources/hibernate.cfg.xml (alla riga 13). 

\begin{xml}
<?xml version="1.0" encoding="utf-8"?>
<!DOCTYPE hibernate-configuration PUBLIC
"-//Hibernate/Hibernate Configuration DTD 3.0//EN"
"http://hibernate.sourceforge.net/hibernate-configuration-3.0.dtd">
<hibernate-configuration>
    <session-factory>
        <property name="hibernate.connection.driver_class">com.mysql.jdbc.Driver</property>
        <property name="hibernate.connection.password"></property>
        <property name="hibernate.connection.url">jdbc:mysql://localhost/clienti</property>
        <property name="hibernate.connection.username">nicola</property>
        <property name="hibernate.dialect">org.hibernate.dialect.MySQLDialect</property>
       
       <mapping resource="com/acme/model/Prodotti.hbm.xml" />
    </session-factory>
</hibernate-configuration>
\end{xml}

Attenzione: il goal hibernate3:hbm2cfgxml cancella e riscrive ogni volta questo file ed inoltre vi aggiunge delle configurazioni che a runtime non sono usate (come username, password, url e driver class). Nell'impiego di tutti i giorni di Regola kit il nostro team di sviluppo non utilizza il goal hibernate3:hbm2cfgxml e si occupa di aggiungere manualmente i mappaggi delle risorse. Naturalmente le configurazioni non usate non costituiscono problema, per cui alla fine la scelta di impiegare o meno hibernate3:hbm2cfgxml è lasciata alla vostra discrezione.

\begin{nota}
Esistono altri goal disponibili, ad esempio la generazione degli script sql per creare le tabella a partire dalla configurazione delle classi. Si rimanda, per approfondimenti, al capitolo \vref{chap:persistenza}.
\end{nota}

\section{Dal modello alla presentazione}
Adesso che avete creato la classe di modello è il momento di realizzare il codice per leggere e scrivere oggetti (della classe Prodotto) sul database, le pagine web che elenchino questi oggetti così come la pagine di dettaglio per effettuare modifiche. Questo codice può essere scritto a mano oppure potete partire facendovi generare automaticamente della classi di default che utilizzerete come modello di partenza per le vostre modifiche.

Per avviare il generatore di Regola kit utilizzate il seguente comando:

\begin{bash}
nicola@casper:~/projects/clienti# mvn exec:java -Dexec.args="-c com.acme.model.Prodotto -m"
\end{bash}

Noterete tra i parametri passati al comando il nome della classe attorno a cui costruire i vari livelli e l'opzione \emph{m} che specifica di utilizzare un'ampia catena di generatori, in particolare:

\begin{center}
{
  \begin{tabular}{ | l | p{9cm} | }
  \hline
  dao &  produce il custom dao \\ \hline
  modelPattern & produce la classe necessaria a Model Pattern   \\ \hline
  properties & aggiunge le chiavi per la localizzazione \\ \hline
  list-handler & genera il controller dietro la pagina di lista \\ \hline
  list &  genera la pagina di lista\\ \hline
  form-handler & genera il controller dietro la pagina di dettaglio\\ \hline
  form &  genera la pagina di dettaglio \\ \hline
  \end{tabular}
}
\end{center}

\begin{nota}
I generatori possono anche essere avviati individualmente. Per scoprire come e conoscere anche altri generatori forniti con Regola kit si rimanda al capitolo \vref{sec:generatori}.
\end{nota}







\chapter{Installazione}\label{chap:installazione}

\section{Maven 2}
TODO: Qui si spiega che Regola kit si fonda su Maven 2 per tutta la gestione del ciclo di build.

\section{Librerie}
TODO: Dove si elencano i moduli di cui si compone

\section{Dipendenze}
TODO: Dove si elencano le dipendenze, ricordando però che sono gestite da Maven 2

\section{Eclipse IDE}
TODO: Come lavorare ad un progetto di Regola utilizzando Eclipse 3.3 ed i plugin necessari per il funzionamento
\chapter{Struttura di un progetto}\label{chap:struttura}
Questo capitolo funziona un po' come una cartina stradale e rimanda ad altri capitoli per l'approfondimento.

\begin{center}
{
  \begin{tabular}{ | l | l | p{6cm} | }
  \hline
  / & pom.xml  & il file principale di Maven 2 \\ \hline

  src/main/resources/ & runtime.properties  & proprietà di runtime \\ \hline
  src/main/resources/ & log4j.xml & la verbosità dei log \\ \hline
  src/main/resources/ & hibernate.cfg.xml & la configurazione principale di Hibernate  \\ \hline
  src/main/resources/ & ApplicationResources.properties & i testi tradotti nelle varie lingue  \\ \hline
  src/main/resources/ & applicationContext-*.xml & le configurazioni di Spring  \\ \hline
  src/main/resources/ & <stesse cartelle dei packages> & i mappaggi di Hibernate  \\ \hline

  src/test/resources/ & designtime.properties  & proprietà a design time \\ \hline
  src/test/resources/ & log4j.xml & la verbosità dei log \\ \hline
  src/main/resources/ &  hibernate.reveng.xml & Hibernate tools  \\ \hline
  src/main/resources/ & applicationContext-*-test.xml & le configurazioni di Spring  \\ \hline
  src/main/resources/ & jetty/env.xml & datasource per Jetty  \\ \hline

  \end{tabular}
}
\end{center}


\section{Lo standard Maven 2}
TODO: Si descrive lo standard utilizzato e si fa una prima panoramica delle cartelle coinvolte

\section{L'iniezione delle dipendenze}
TODO: La distinzione tra i tre livelli di bean

\section{La localizzazione}
TODO: I file di localizzazione

\section{Le connessioni al database}
TODO: I file di localizzazione


\section{Verbosità dei log}
TODO: log4j e come interagisce con JBoss

\section{La sezione dei test}
TODO: come configurare i test e quali file utilizzano.

\section{Applicazione Servlet}
TODO: Davvero in breve la struttura
\chapter{Database}\label{chap:database}\index{datasource|see{database}}

\section{Configurazione di run-time}\index{database!run-time config}
Le applicazioni JEE generalmente lasciano la gestione delle connessioni database al container. In questo modo è possibile per i sistemisti modificare, ad esempio, la url di connessione oppure il numero di connessioni in pool senza modificare l'applicazione né dovere effettuare una riconsegna (redeploy). Ogni container configura in modo diverso però ogni datasource è caratterizzato necessariamenta da un nome JDNI, ad esempio java:comp/env/jdbc/miodatabase. Questo nome è utilizzato dall'applicazione per recuperare la connessione e, in definitiva, connettersi al database. 
Nelle applicazioni Regola kit il nome JNDI del datasource è specificato nella proprietà di runtime jee.datasrouce (nel file  src/main/resource/runtime.properties). 

\begin{bash}
jee.datasource=java:comp/env/jdbc/Datasource
\end{bash}

Con questo la configurazione di runtime della vostra applicazione è terminata, non ci sono altre variabili da impostare. Un problema tipico di configurazione riportato di diversi utenti è l'impossibilità di connettersi al datasource per avere utilizzato un indirizzo JNDI sbagliato. La causa del problema  è che il nome con cui l'application server espone la connessione non è \emph{esattamente} quello specificato nella configurazione: perché ad esempio viene preposta la stringa `java:` o a volte `java:comp/env/`.  
\\\index{database!esempi!JBoss}
Per facilitare la soluzione di questo tipo di problema concludiamo il paragrafo riportando qualche configurazione di datasource per gli application server più diffusi. Ad esempio  JBoss richiede di ricopiare nella cartella di deploy del server utilizzato un file con il nome del tipo *-ds.xml (ad esempio miadatasource-ds.xml) con un contenuto simile al seguente:

\begin{xml}
<?xml version="1.0" encoding="UTF-8"?>

<datasources>
 <local-tx-datasource>
    <jndi-name>jdbc/services</jndi-name> 
    
    <connection-url>jdbc:oracle:thin:@133.222.0.1:1522:SID</connection-url>
    <user-name>username</user-name>
    <password>*****</password>
  
    <driver-class>oracle.jdbc.driver.OracleDriver</driver-class>
    <exception-sorter-class-name>
       org.jboss.resource.adapter.jdbc.vendor.OracleExceptionSorter
    </exception-sorter-class-name>
    <min-pool-size>5</min-pool-size>
    <max-pool-size>20</max-pool-size>

 </local-tx-datasource>
</datasources>
\end{xml}

In questo caso nella configurazione dell'applicazione dovete indicare il nome indicato però preceduto da `java:` come indicato di seguito.

\begin{bash}
jee.datasource=java:jdbc/services
\end{bash}


\index{database!esempi!Jetty}Per quanto riguarda Jetty il file di configurazione env.xml contiene l'indicazione del datasource.

\begin{xml}
<?xml version="1.0"?>
<!DOCTYPE Configure PUBLIC "-//Mort Bay Consulting//DTD Configure//EN" "http://jetty.mortbay.org/configure.dtd">
<Configure class="org.mortbay.jetty.webapp.WebAppContext">

  <New id="jira-ds" class="org.mortbay.jetty.plus.naming.Resource">
     <Arg>jdbc/Datasource</Arg>
       <Arg>
         <New class="org.enhydra.jdbc.standard.StandardConnectionPoolDataSource">
           <Set name="Url">jdbc:mysql://localhost/clienti</Set>
           <Set name="DriverName">com.mysql.jdbc.Driver</Set>
           <Set name="User">nicola</Set>
         </New>
      </Arg>
  </New>

</Configure>

\end{xml}

Jetty aggiunge al nome configurato la stringa `java:comp/env/` per cui l'indirizzo JNDI da utilizzare è il seguente:

\begin{bash}
jee.datasource=java:comp/env/jdbc/Datasource
\end{bash}

TODO: Aggiungere la configurazione per Tomcat

TODO: Mi chiedo come mai il dialetto sia specificato come proprietà di designtime ma non come proprietà di runtime.



\section{Configurazione di design-time}\index{database!design-time config}
Tra le caratteristiche invidiabili del framework Spring vi è la possibilità di testare l'applicazione fuori dal container con notevoli risparmio di tempo e relativo aumento di produttività. I servizi necessari ai test sono offerti direttamente da Spring che si sostituisce al container, ad esempio, nella gestione delle transazioni e nei datasource. In quest'ultimo caso è necessario specificare tutte le proprietà di una connessione (come il driver, la url, la username e la password) come proprietà di designtime (nel file src/test/resources/designtime.properties).

\begin{bash}
...
hibernate.dialect=org.hibernate.dialect.MySQLDialect
hibernate.connection.driver_class = com.mysql.jdbc.Driver
hibernate.connection.url = jdbc:mysql://localhost/clienti
hibernate.connection.username = nicola
hibernate.connection.password = 
...
\end{bash}

Da notare che è necessario specificare anche la variabile hibernate.dialect con il tipo di \emph{dialetto} di database utilizzato. I tipi disponibili sono molti ed elencati nella documentazione di Hibernate, alcuni esempi comuni sono MySQL5Dialect, OracleDialect, PostgreSQLDialect, HSQLDialect e SQLServerDialect.




\chapter{Persistenza}\label{chap:persistenza}

\section{Hibernate}
TODO: Qui si dice che l'orm predefinito da Regola kit è Hibernate, perché è estremamente flessibile e su database preesistenti consente di gestire anche i casi più anomali. Perché non usiamo ancora JPA (immaturo).

\section{Configurazione}
TODO: I file hibernate.cfg.xml ed i mappaggi. Alcuni esempi di base per questi file.

\section{Generatori automatici}
TODO: Esempi di configurazione di hibernate.reveng.xml ed esempi (tutti) dei vari goals disponibili a riga di comando.

\section{Altri ORM}
TODO: Si parla del supporto sperimentale per gli altri orm.
\chapter{Messa in produzione}\label{chap:produzione}

\section{Produrre i pacchetti war ed ear}
TODO: Qui si spiega come produrre i pacchetti e cosa contengano al loro interno

\section{Application Server}\label{sec:datasource}
TODO: Perché non usare Jetty ed esempi di configurazione per Tomcat e JBoss. Si rimanda al capitolo dei database per la configurazione dei datasource.

Per utilizzare CXF all'intenro di JBOSS è necessario impostare il parametro di avvio della virtual machine javax.xml.soap.MessageFactory al valore com.sun.xml.messaging.saaj.soap.ver1\_1.SOAPMessageFactory1\_1Impl.



\section{Integrazione continua}
TODO: Perché serve e come configurare i file standard dei progetti di Regola kit
\chapter{Sviluppo}
Questo capitolo apre la parte dedicata alle funzionalità offerte da Regola kit relativamente allo sviluppo. Lasciandosi alle spalle le configurazioni.

\section{Domain Driven Development}
TODO: Molto in breve si spiega come ci si incentri sul modello.

\section{Livelli}
TODO: Si presentano i vari livelli e si rimanda a capitoli successivi per i dettagli

\section{Model Pattern}\label{sec:modelpatternTeoria}
\index{filtri|see{model pattern}}
\index{selezione|see{model pattern}}
\index{model pattern!presentazione}
\index{estrazione|see{model pattern}}
Regola kit nasce per agevolare la realizzazione di applicazioni web destinate a gestire moli piuttosto ampie di dati generalmente persistite su un database relazionale. Un problema ricorrente in questi contesti è quello che si potrebbe sinteticamente indicare come  come la questione dei sottoinsiemi, ovvero fornire una risposta generale alle domande:  come estrarre in modo conveniente sottoinsiemi di una certa entità dal database? E come rappresentare in modo sintetico questi sottoinsiemi all'interno di un form html oppure in un'applicazione desktop?
\\
La soluzione proposta da Regola kit  riteniamo che possa essere ritenuta sufficientemente assodata e generica da ambire al rango di design pattern (magari un minore); se così fosse pensiamo che la descrizione del problema unitamente alla soluzione individuata possa essere identificata col nome di Model Pattern (o in alternativa Selection).\\
 
Nei prossimi paragrafi cercheremo di formalizzare questo nuovo pattern descrivendone intenti e forze in gioco. Si avvisa fin d'ora che materremo il discorso a livello teorico presentando solo esempi astratti utili unicamente per comprendere la portata della soluzione proposta. Rimandiamo al paragrafo \vref{sec:modelpattern} per illustrare l'implementazione di riferimento di Model Pattern (quella contenuta in Regola kit) e spiegarne l'utilizzo all'interno delle vostre applicazioni. 

\subsection{Intento}
Model Pattern vuole fornire un design appropriato per individuare un sottoinsieme di oggetti del modello di dominio e rappresentare questo sottoinsieme in modo conveniente all'interno del livello di presentazione.

\subsection{Forze}
In una categoria piuttosto ampia di applicazioni (che comprende la quasi totalità di applicazioni web)  gli oggetti del modello salvano il loro stato in modo persistente su database, magari attraverso dei framework ORM (ad esempio Hibernate, JPA, ecc.). Non di rado il numero di istanze persistenti, per tipologia di classe, è veramente elevato cosa che complica in primis l'estrazione (1) (ovvero il filtraggio) delle istanze. Sempre dalla mole dei dati deriva la necessità di ordinare (2) e paginare (3) gli oggetti estratti: ovvero raggrupparli in sotto gruppi contenenti al più un certo numero di elementi.
Un altra richiesta tipica delle applicazioni di cui stiamo parlando è quella di rappresentare gli oggetti del modello estratti in modo parziale, ovvero mostrare solo un sotto gruppo di proprietà, in modo da fornire, in un solo colpo d'occhio, tutti gli elementi necessari per prendere delle decisioni sui dati eliminando le proprietà superflue. Questo sezionamento delle classi del modello si chiama proiezione (4).\\
Del sottoinsieme degli elementi del modello individuato e estratto dal database (reidratato) è inoltre necessario fornire delle rappresentazioni sintetiche da utilizzare in diversi punti dell'applicazione. Si prenda per esempio la classe Studente ed  il form html che consente di effetture le operazione di editing sulle sue proprietà quali nome, cognome ed altre. Si consideri poi che tra le proprietà della classe Studente vi sia anche la collezione di classi del tipo Esami che contenga gli esami sostenuti fino alla data corrente. Bene, il problema è spesso quello di mostrare nello stesso form html contenete le proprietà di Studente anche la collezione di Esami: in questo caso la rappresentazione corretta potrebbe variare in base alle finalità applicative. Ad esempio le segreterie di facoltà potrebbero essere interessati all'elenco completo di tutti gli esami (nessuna sintesi), i docenti potrebbero volere esclusivamente la media delle valutazione degli esami (applicare una metrica), l'ufficio delle imposte potrebbe essere interessato solo al fatto di aver sostenuto almeno un esame (e quindi un intero che esprime il numero di esami superati). Tutte queste rappresentazioni (elenco, metrica ed intero) sono una descrizione del sottoinsieme di oggetti del modello Esami.\\
Inoltre si consideri il caso in cui la collezione di Esami fosse vuota o nulla, in questo caso diverse rappresentazioni tradurrebbero il valore null nel modo più appropriato rispetto al contesto (ad esempio con una stringa che indica l'assenza di esami, o con una segnalazione di errore o quant'altro).  In ogni caso per ogni sottoinsieme di oggetti del modello emerge prepotente la necessità di provvedere ad una o più rappresentazione sintetica (5).\\
Quindi ricapitolando le forze che caratterizzano il problema:
\begin{enumerate*}
  \item estrazione di dati efficiente da una mole ampia o ampissima
  \item ordinare  le istanze di modello estratte
  \item paginare  le istanze di modello estratte
  \item effettuare delle proiezioni sulle istanze di modello estratte
  \item fornire delle rappresentazioni sintetiche della totalità di istanze estratte (cioè del sottoinsieme di tutta la popolazione di istanze estratto)
\end{enumerate*}

\subsection{Esempio}
Per amore di concretezza riportiamo un esempio che possa aiutare a comprendere meglio il problema descritto nel paragrafo precedente e la soluzione proposta da Model Pattern. Immaginiamo di avere una collezione di oggetti persista su database, tutti appartenenti alla medesima classe Prodotto, che presentano solo due proprietà: id, un intero che tiene la chiave primaria, e la descrizione del prodotto stesso. 

\begin{java}
public class Prodotto {
    
    private Integer id;
    private String descrizione;

    //getter and setter per tutti i campi
    ...
}
\end{java}

Il problema è quello di descrivere in modo sintetico sottoinsiemi di oggetti della classe Prodotto. Procederemo per gradi iniziando con il sottoinsieme diciamo $\Sigma$ definito in modo sintetico, dalla descrizione seguente: il sottoinsieme di tutti gli oggetti la cui descrizione inizia con la parola 'manuale'. L'equivalente informatico di questa definizione è costiuito da un'instanza della classe seguente (che svolge il ruolo di Model Pattern):

\begin{java}
public class ProdottoPattern  {
    
    private String inizioDescrizione;

    //getter and setter per tutti i campi
    ...
}
\end{java}

In particolare il sottoinsieme $\Sigma$ può essere definito da un oggetto della classe ProdottoPattern che ha come proprietà inizioDescrizione la stringa 'manuale'. Si può notare che ProdottoPattern non contiene né il sottoinsieme $\Sigma$ né il codice per estrarlo da database; si limita solo a fornirne una descrizione sintetica del sottoinsieme.
Arricchiamo questo primo esempio includendo anche l'ordine degli oggetti contenuti nel sottoinsieme $\Sigma$, ad esempio richiedendo che siano ordinati in base alla proprietà id. In questo caso la classe Model Pattern potrebbe modificarsi così:

\begin{java}
public class ProdottoPattern  {
    
    private String inizioDescrizione = ``manuale'';
    private String[] ordine = { ``id'' }; 

    //getter and setter per tutti i campi
    ...
}
\end{java}

In questo modo sarebbe in grado di descrivere il sottoinsieme ordinato $\Sigma$ anche se in modo un po' rozzo infatti si trascura la possibilità di ordinare in modo ascendente o discendente. Comunque l'esempio dovrebbe rendere l'idea di come le classi Model Pattern si limitino a descrivere i sottoinsiemi senza contenere codice per l'estrazione dello stesso. Questo servizio è infatti fornito dalle classi a corredo come quelle fornite da Regola kit che consentono, dato un oggetto del tipo Model Pattern, di ottenere il sottoinsieme voluto. La sfida consiste nella predisposizione di un architettura che consenta di utilizzare dei Model Pattern semplici come quelli mostrati in questi esempi ma in grado di funzionare correttamente (ad esempio indicando che la proprietà inizioDescrizione si riferisce alla proprietà descrizione di Prodotto) su diversi orm. Si rimanda al paragrafo successivo per una descrizione esatta dell'architettura.\\
Fino a questo momento abbiamo trattato solo la capacità di Model Pattern di esprimere una selezione, adesso ci occuperemo della funzione di rappresentazione del sottoinsieme selezionato. Immaginiamo di dover mostrare in una pagina web il sottoinsieme $\Sigma$ in modo sintetico (ovvero senza elencare in una colonna tutti i suoi elementi). Modifichiamo Model Pattern per aggiungere la proprietà sintesi.

\begin{java}
public class ProdottoPattern  {
    
    public getSintesi()
    {
       if (inizioDescrizione == null)
       {
          return ``Prodotti con descrizione iniziante con: '' +  inizioDescrizione;
       }
       else
       {
          return "Tutti di prodotti.";
       }
    }

}
\end{java} 

Come si vede la proprietà sintesi contiene la logica per fornire una semplice rappresentazione del sottoinsieme $\Sigma$ anche nel caso in cui inizioDescrizione sia nullo.



\subsection{Architettura}
Model Pattern propone di avvilire le forze del problema utilizzando una classe (Pattern) per ogni tipo nel modello che si occupi di due distinti aspetti:
\begin{enumerate*}
  \item individuare il sottoinsieme selezionato
  \item fornire (zero o più) rappresentazione del sottoinsieme selezionato
\end{enumerate*}

Per quanto riguarda il primo punto bisogna precisare fin da subito che la classe Pattern non contiene in nessun modo il codice necessario all'estrazione della selezione o la selezione stessa. Invece si occupa di contenere le informazioni necessarie a descrivere in modo esatto il sottoinsieme,  ovvero descriverlo in modo esatto rispetto ai seguenti punti:

\begin{enumerate*}
\item quali istanze sono comprese nel sottoinsieme e quali esluse
\item con quali ordine le istanze entrano nel sottoinsieme
\item in che modo il sottoinsieme è paginato
\item quali proiezioni effettuare su ogni singola istanza del sottoinsieme
\end{enumerate*}

La classe Pattern si occupa anche di contenere zero o più rappresentazioni del sottoinsieme che descrive da utilizzare in altrettanti punti dell'applicazione (tipicamente un qualche livello di presentazione, o la produzione di report, ecc.). A tal proposito predispone  il metodo

\begin{java}
    void init(Collection<Model> subset) 
\end{java}

questo metodo prende in ingresso un sottoinsieme e provvede ad inizializzare le rappresentazioni di subset secondo la logica contenuta in pattern; 

[TODO: nda Lorenzo, in futuro potrebbe essere utile utilizzare questo stesso metodo per inizializzare anche la parte di descrizione 1) in base ad un subset, ovvero produrre il filtro che estrarrebbe eventualmente subset?]


\section{Generatori}\label{sec:generatori}
TODO: Qui si elencano i generatori disponibili e cosa scrivano. Se il paragrafo diventasse troppo lungo allora lo mettiamo su un capitolo a parte.
\chapter{Dao}

\section{Scopo}
TODO: A cosa serve DAO?

\section{GenericDao}
TODO: Se ne descrivono interfacce e si presentano esempi d'uso (dei test d'unità) che si snodano tra i vari paragrafi.

\section{Creare un custom dao}
TODO: l'interfaccia, l'implementazione ed infine la configurazione di Spring. Si ricorda che esiste un generatore per questo.

\section{Ricerche con Model Pattern}\label{sec:modelpattern}
Nel paragrafo \vref{sec:modelpatternTeoria} è stato presentato in modo formale un nuovo design pattern chiamato Model Pattern; è stato spiegato come intenda risolvere il probema dell'estrazione di un sottoinsieme di oggetti e di come intenda fornire una rappresentazione sintetica di questo oggetto. La soluzione proposta, si è visto, prevede l'utilizzo di una classe che svolga il ruolo di ModelPattern, ovvero sia contemporaneamente in grado di individuare quali elementi estrarre e di rappresentarli senza però contenere al suo interno né il sottoinsieme né la logica per l'estrazione. La discussione è rimasta però a livello astratto rimandando la descrizione di come Model Pattern possa essere utilizzato concretamente all'interno di un'applicazione Regola kit; nei prossimi paragrafi vedremo quindi l'implementazione di riferimento di Model Pattern contenuta in Regola kit (precisamente nei moduli regola-core, regola-dao e sottomoduli) per scoprire come progettare un classe ModelPattern con diversi criteri di filtraggio e ordinamento e la utilizzeremo lungo i diversi livelli applicativi, dal DAO alla presentazione. Per maggiore concretezza immagineremo di doverci occupare, diciamo, della classe Prodotto riportata di seguito.

\begin{java}
class Prodotto {
    Integer id;
    String  descrizione;

    //seguono getter/setter per ogni campo
    ...
}
\end{java}

\subsection{ModelPattern}
Regola kit fornisce una classe di base org.regola.model.ModelPattern da cui derivare il nostro ModelPattern; questa classe presenta alcune facilitazioni per specificare gli ordinamenti e la paginazione ed è accettata quasi in ogni livello applicativo, dai GenericDao fino a controllori web che si occupano di disegnare la tabella con il sottoinsieme. Tenendo a mente la classe Prodotto una prima versione del nostro ModelPattern potrebbe essere la seguente:

\begin{java}
class ProdottoPattern extends org.regola.model.ModelPattern 
    implements Serializable {
    
    Integer  chiave;

    @Equals("id")
    public Integer getChiave() {
        return chiave;
    }

    ...
}
\end{java}

Soffermiamoci un attimo su alcuni elementi che trasformano la nostra classe in un ModelPattern utilizzabile con Regola kit:

\begin{description*}
  \item[nome della classe] una convenzione di Regola kit prevede che ogni ModelPatter dovrebbe chiamarsi col nome della classe di modello a cui si riferisce (nel nostro caso Prodotto) seguito dal suffisso Pattern, per cui ProdottoPattern
  \item[derivazione] la nostra classe deve derivare dalla classe org.regola.model.ModelPattern fornita con Regola kit
\item[serilizzazione] la nostra classe deve essere progettata per la serializzazione e quindi necessariamente implementare l'interfaccia di marker java.io.Serializable
\end{description*}

La classe ProdottoPattern così realizzata è predisposta per individuare particolari sottoinsiemi di instanze Prodotto che presentano particolari valori della proprietà id, in particolare tutte le istanze di Prodotto che hanno id uguale alla proprietà chiave di ProdottoPattern. Il legame tra chiave ed id è realizzato mediante l'annotazione @Equals (alla riga 6) che, posta sul getter della proprietà chiave di ProdottoPattern, unisce quest'ultima alla proprietà dell'oggetto di modello specificata dentro l'annotazione stessa (in questo caso id).
\\
Bisogna chiarire da subito che non è la classe ModelPattern ad individuare un sottoinsieme ma ogni sua istanza per cui, tornando al nostro esempio, bisogna istanziare un oggetto del tipo ProdottoPattern. 

\begin{java}
ProdottoPattern pattern = new ProdottoPattern();

pattern.setChiave(234532); //ora pattern rappresenta un sottoinsieme
\end{java}

Dopo l'assegnazione pattern (l'oggetto) e non ProdottoPattern (la classe) è in grado di individuare un sottoinsieme di oggetti Prodotto aventi la proprietà id uguale a 234532 (essendo id una chiave primaria il sottoinsieme conterrà al più un elemento). Da notare come l'oggetto pattern non contenga in sé il sottoinsieme ma solo la descrizione sintetica di questo; per ottenere il sottoinsieme bisogna rivolgersi alla classe ProdottoDao invocando il metodo find.

\begin{java}
ProdottoPattern pattern = new ProdottoPattern();

pattern.setChiave(234532); 

List<Prodotto> prodotti = prodottoDao.find(pattern);
\end{java}

Ora che abbiamo visto concretamente come creare ed utilizzare una classe ModelPattern possiamo scendere nel dettaglio ricordando che Model Pattern deve individuare con esattezza il sottoinsieme, in particolare deve occuparsi di ordinamento, paginazione e proiezione. Nei prossimi paragrafi quindi scopriremo come specificare:

\begin{enumerate*}
  \item la selezione da effettuare
  \item i criteri di ordinarmento
  \item la paginazione da effettuare 
  \item le proprietà da comprendere nella proiezione
  \item fornire rappresentazioni sintetiche del sottoinsieme indicato
\end{enumerate*}

\subsection{La selezione}
La selezione avviene impostando alcuni criteri di filtraggio sul ModelPattern. Ogni criterio è caratterizzato da:

\begin{enumerate*}
  \item la proprietà del modello a cui si applica
  \item il tipo di criterio (ad esempio ugualianza, appartenenza, ecc.) 
  \item il valore da utilizzare per il confronto
\end{enumerate*}

Riprendendo l'esempio del paragrafo precedente troviamo un unico criterio di filtraggio, di tipo ugualianza, che confronta la proprietà di modello id con il valore della proprietà chiave di ProdottoPattern.

\begin{java}
class ProdottoPattern extends org.regola.model.ModelPattern 
    implements Serializable {
    
    Integer  chiave;

    @Equals("id")
    public Integer getChiave() {
        return chiave;
    }

    ...
}
\end{java}

La proprietà di modello può trovarsi direttamente sulla classe radice, ovvero quella a cui riferisce ModelPattern (nell'esempio Prodotto) oppure su classi a questa collegate tramite associazioni del tipo uno-a-uno, molti-a-uno od uno-a-molti. In generale si navigano le associazioni utilizzando come separatore il punto tranne che per relazioni uno-a-molti (collezioni) dove si utilizza il simbolo [] postfisso. Nel dettaglio:

\begin{center}
{
  \begin{tabular}{ | l | l | l  | }
  \hline
  radice  & il nome della proprietà & id \\   
                                    & & nome \\   \hline
  uno-a-uno  & il punto    & indirizzo.via \\   
                         & & categoria.descrizione \\   \hline
  molti-a-uno  & il punto    & cliente.nome \\   
                           & & fattura.progressivo \\   \hline
  uno-a-molti  & [] postfisso ed il punto    & elmenti[].nome \\   
                                           & & elementi[].dettagli[].progressivo \\   \hline
  \end{tabular}
}
\end{center}

Il tipo del criterio è impostato scegliendo l'annotazione tra quelle presenti in Regola kit. Attualmente è possibile scegliere tra le annotazioni seguenti:

\begin{center}
{
  \begin{tabular}{ | l | l  | }
  \hline
    Equals & ugualianza \\ \hline
    NotEquals & disugualianza \\ \hline
    Like & il like \\ \hline
    GreatherThan & maggiore \\ \hline 
    LessThan & minore \\ \hline
    In & appartenenza \\ \hline
  \end{tabular}
}
\end{center}

Il valore di ogni criterio una proprietà di ModelPattern individuata annotandone il getter. Il tipo di questa proprietà deve corrispondere a quello del modello tranne per il criterio In per cui è necessario specificare un'array. L'esempio seguente fornisce una rappresentazione piuttosto completa di quando esposto:

\begin{java}
public class CustomerPattern extends ModelPattern 
   implements Serializable {
  
  private Integer id;

  @Equals("id")
  public .Integer getId()
  {
    return id;
  }
  
  private String firstName;
  
  @Like(value = "firstName", caseSensitive = true)
  public String getFirstName()
  {
    return firstName;
  }
    
  private String[] lastNames;
  
  @In("lastName")
  public String[] getLastNames() {
    return lastNames;
  }

  private Integer invoiceId;
 
  @Equals("invoices[].id")
  public Integer getInvoiceId() {
    return invoiceId;
  }

}
\end{java}

\subsection{Ordinamento}
Sugli oggetti selezionati è possibile impostare criteri di ordinamento in base alle proprietà della radice del modello o sugli oggetti ad essa associati. L'implementazione di Model Pattern presente in Regola kit utilizza per specificare gli ordinamenti (così come le proiezioni) la classe ModelProperty, una specie di descrittore di una generica proprietà. Si instazia così:

\begin{java}
   ModelProperty mp =  new ModelProperty("id","customer.column.",Order.asc);   
\end{java}

ModelProperty contiene il nome della proprietà (in base a cui ordinare), un prefisso da utilizzare per individuare in modo univoco la proprietà all'interno di tutta l'applicazione ed, infine, la direzione dell'ordinamento (acendente o discendente). 
Nell'esempio si esprime un ordinamento ascentente sulla proprietà id, individuata nell'applicazione come "customer.column.id".
\\
Per speficiare l'ordinamento basta popolare la lista sortedProperties di ModelPattern con tante istanze di ModelProperty in base all'ordinamento da realizzare. Ad esempio:

\begin{java}
   ModelProperty id =  new ModelProperty("id","customer.column.",Order.asc);
   ModelProperty street =  new ModelProperty("address.street","customer.column.",Order.desc);

   modelPattern.getSortedProperties.clear();
   modelPattern.getSortedProperties.add(id);
   modelPattern.getSortedProperties.add(street);
\end{java}

In questo caso si impone un ordinamento crescente per id e descrescente per la proprietà strret dell'oggetto associato address.


\subsection{Paginazione}
Sugli oggetti così selezionati ed ordinati è possibile effettuare un'olteriore selezione dividendolo in blocchi contigui dette pagine contenenti al più $n$ elementi. ModelPattern consente di impostare la dimensione delle pagine così come impostare la pagina da selezionare. Ad esempio:

\begin{java}
  modelPattern.setPageSize(20);
  modelPattern.setCurrentPage(0);
\end{java}

Qui si suddivide l'insieme ordinato in pagina con al più di 20 oggetti, il primo blocco comprende gli oggetti da $[0,19]$, il secondo da $[20,39]$ e così via. Inoltre si seleziona la prima pagina (la numerazione delle pagine parte da 0), ovvero gli oggetti da $[0,19]$.

\subsection{Proiezione}
Fino a qui abbiamo visto come effettuare la selezione, ordinarla e limitarla ad un certo numero di elementi. \'{E} inoltre possibile limitare non solo il numero ma anche le proprietà del singolo oggetto del sottoinisieme. Ad esempio può essere conveniente limitarsi a considerare solo la proprietà id e name piuttosto che il complesso di tutte le proprietà della radice del modello e/o degli oggetti associati. Questo genere di limitazione si chiama proiezione.
Il meccanismo per specificare le proiezione dentro Regola kit si basa sempre sugli oggetti ModelProperty, basta popolare la lista visibleProperties di ModelPattern con le proprietà che si intende proiettare.


\begin{java}
   ModelProperty id =  new ModelProperty("id","customer.column.",Order.asc);
   ModelProperty street =  new ModelProperty("address.street","customer.column.",Order.desc);

   modelPattern.getVisibleProperties.clear();
   modelPattern.getVisibleProperties.add(id);
   modelPattern.getVisibleProperties.add(street);
\end{java}

Nell'esempio si effettua una proiezione comprendente esclusivamente le proprietà id e strett dell'oggetto associato street.

\subsection{Rappresentazione}
TODO:























\chapter{Servizio}

\section{Scopo}
TODO: A cosa serve il livello di Servizio?

\section{GenericManager}
TODO: Come crearlo. Si ricorda che esiste un generatore per questo.

\section{Transazioni}
TODO: come sono demarcate e come aggiungere politiche diverse per aprire e chiudere transazioni

\section{Politiche di detach}
TODO: come usare oggetti detacciati, la sessione lunga, il lazy loading ecc.

\section{Web Services}
TODO: esempi di configurazioni già pronte dentro Regola kit
\chapter{Presentazione Web}

\section{Scopo}
TODO: Di cosa si occupa questo livello?

\section{Tecnologie}
TODO: Panoramica brevissima su JSF, Spring WebFlow e Spring MVC.

\section{Pagina di lista}
TODO: Partendo da un esempio si provvede a spiegare quali sono i file coinvolti e come devono essere configurati. Si ricorda che esiste un generatore per questo.

\section{Pagina di form}
TODO: Partendo da un esempio si provvede a spiegare quali sono i file coinvolti e come devono essere configurati. Si ricorda che esiste un generatore per questo.


% -----------BIBLIOGRAFIA

%\begin{thebibliography}{9}
%\bibitem{lamport94}
%  Leslie Lamport,
%  \emph{\LaTeX: A Document Preparation System}.
%  Addison Wesley, Massachusetts,
%  2nd Edition,
%  1994.
%\end{thebibliography}

\bibliographystyle{alpha}
\bibliography{main}



% -----------INDICE
\printindex

\end{document}
